\setcounter{section}{22}

  \zwischenueberschrift{Soal latihan}

\inputaufgabe
{}
{

Misalkan $M$ suatu
\definitionsverweis {manifold berbatas}{}{}
dan
\mathl{\partial M}{} merupakan batasnya. Tunjukkan bahwa
\definitionsverweis {batas topologis}{}{}
dari
\mathl{M \setminus \partial M}{} sama dengan $\partial M$.

}
{} {}

\inputaufgabe
{}
{

Tentukan
\definitionsverweis {tumpuan}{}{} fungsi-fungsi berikut dari $\R$ ke $\R$.
\aufzaehlungsieben{Suatu fungsi polinomial.
}{Fungsi sinus.
}{Fungsi eksponensial.
}{Fungsi indikator $e_{ \Z }$.
}{Fungsi indikator $e_{ \Q }$.
}{Fungsi indikator $e_{  [a,b]  }$.
}{Fungsi indikator $e_{ ]a,b[ }$.
}

}
{} {}

\inputaufgabe
{}
{

Misalkan
\maabbdisp {f} { X } {\R
} {}
suatu fungsi pada
\definitionsverweis {ruang topologis}{}{}
$X$. Andaikan terdapat suatu
\definitionsverweis {himpunan bagian terbuka}{}{}

\mavergleichskette
{\vergleichskette
{U
}
{ \subseteq }{ X
}
{  }{
}
{    }{
}
{   }{
}
}
{}{}{,}
yang memuat
\definitionsverweis {tumpuan}{}{}
dari $f$. Tunjukkan bahwa $f$
\definitionsverweis {kontinu}{}{}
jika dan hanya jika pembatasan
\mathl{f \vert_U}{} kontinu.

}
{} {}

\inputaufgabe
{}
{

Misalkan $X$ suatu
\definitionsverweis {ruang topologis}{}{}
dan

\mavergleichskette
{\vergleichskette
{ T
}
{ \subseteq }{ X
}
{  }{
}
{    }{
}
{   }{
}
}
{}{}{}
suatu himpunan bagian. Tunjukkan bahwa
\definitionsverweis {penutupan}{}{}
dari $T$ sama dengan
\definitionsverweis {tumpuan}{}{}
dari
\definitionsverweis {fungsi indikator}{}{}

\mathl{e_{  T  }}{}.

}
{} {}

\inputaufgabe
{}
{

Berikan suatu
\definitionsverweis {penghabisan kompak}{}{} untuk
\definitionsverweis {bilangan riil}{}{} $\R$.

}
{} {}

\inputaufgabegibtloesung
{}
{

Berikan suatu
\definitionsverweis {penghabisan kompak}{}{}
untuk $\R^d$.

}
{} {}

\inputaufgabe
{}
{

Misalkan $X$ suatu
\definitionsverweis {ruang topologis}{}{.}
Untuk tutupan atas $X$ yang hanya terdiri atas $X$, tentukan suatu
\definitionsverweis {partisi satuan subordinat terhadap tutupan tersebut}{}{.}

}
{} {}

\inputaufgabe
{}
{

Misalkan $X$ suatu
\definitionsverweis {ruang topologis}{}{}
dan

\mavergleichskette
{\vergleichskette
{ X
}
{ = }{ \bigcup_{i \in I} U_i
}
{  }{
}
{    }{
}
{   }{
}
}
{}{}{}
suatu
\definitionsverweis {tutupan terbuka}{}{.}
Kita tinjau keluarga
\definitionsverweis {fungsi indikator}{}{}

\mathbeddisp {e_{  P  }} {}
 {P \in X} {}
 {} {} {} {.}
Untuk tutupan yang dimaksud
\zusatzklammer {tutupan ini} {}{}{,}
manakah sifat-sifat suatu
\definitionsverweis {partisi satuan subordinat}{}{}
terhadap tutupan tersebut yang dipenuhi oleh keluarga ini?

}
{} {}

\inputaufgabe
{}
{

Kita tinjau
\definitionsverweis {penghabisan kompak}{}{}

\mathbed {A_n= [-n,n]} {}
 {n \in \N} {}
 {} {} {} {,} dari
\definitionsverweis {bilangan riil}{}{}
beserta tutupan terbuka

\mathbed {W_n= A_{n+1}^{o}  \setminus A_{n-1}} {}
 {n \in \N} {}
 {} {} {} {,}
\zusatzklammer {dengan
\mathl{A_{-1}= \emptyset}{}} {} {.}
Carilah suatu tutupan $\R$ dengan
\definitionsverweis {interval terbuka}{}{,}
yang memenuhi sifat-sifat
dari Lemma 22.9
\zusatzklammer {beserta pembuktiannya} {} {}.

}
{} {}

\inputaufgabe
{}
{

Tunjukkan bahwa tidak terdapat barisan
\definitionsverweis {fungsi kontinu}{}{}
\maabbdisp {f_n} { [0,1] } { [0,1]
} {}
untuk

\mavergleichskette
{\vergleichskette
{ n
}
{ \in }{ \N
}
{  }{
}
{    }{
}
{   }{
}
}
{}{}{,}
yang membentuk suatu
\definitionsverweis {partisi satuan}{}{}
dan memiliki sifat bahwa $0$ termasuk dalam
\definitionsverweis {tumpuan}{}{}
setiap fungsi $f_n$.

}
{} {}

\inputaufgabe
{}
{

Misalkan

\mavergleichskette
{\vergleichskette
{ V
}
{ \subseteq }{ H
}
{  }{
}
{    }{
}
{   }{
}
}
{}{}{}
suatu
\definitionsverweis {himpunan terbuka}{}{}
di dalam
\definitionsverweis {setengah ruang}{}{}

\mavergleichskette
{\vergleichskette
{ H
}
{ \subset }{ \R^n
}
{  }{
}
{    }{
}
{   }{
}
}
{}{}{}
dan misalkan
\maabbdisp {f} { V } {\R
} {}
suatu
\definitionsverweis {fungsi yang terdiferensialkan secara kontinu}{}{.}
Tunjukkan bahwa terdapat suatu himpunan terbuka

\mavergleichskette
{\vergleichskette
{ \tilde{V}
}
{ \subseteq }{ \R^n
}
{  }{
}
{    }{
}
{   }{
}
}
{}{}{}
dan suatu
\definitionsverweis {fungsi yang terdiferensialkan secara kontinu}{}{}
\maabbdisp {\tilde{f}} {\tilde{V}} { \R
} {}
sedemikian sehingga

\mavergleichskette
{\vergleichskette
{ \tilde{V} \cap H
}
{ = }{ V
}
{  }{
}
{    }{
}
{   }{
}
}
{}{}{}
dan

\mavergleichskette
{\vergleichskette
{ \tilde{f} \vert_V
}
{ = }{ f
}
{  }{
}
{    }{
}
{   }{
}
}
{}{}{}
berlaku.

}
{} {}

\inputaufgabe
{}
{

Misalkan
\maabbdisp {f} {\R} {\R
} {}
suatu
\definitionsverweis {fungsi yang terdiferensialkan secara kontinu}{}{}
dengan
\definitionsverweis {tumpuan kompak}{}{.}
Tunjukkan bahwa
\definitionsverweis {turunan}{}{}
$f'$ juga memiliki tumpuan kompak dan bahwa

\mavergleichskette
{\vergleichskette
{ \int_\R f' d \lambda^1
}
{ = }{ 0
}
{  }{
}
{    }{
}
{   }{
}
}
{}{}{}
berlaku.

}
{} {}

\inputaufgabe
{}
{

Misalkan
\maabbdisp {f} {\R_{\geq 0}} {\R
} {}
suatu
\definitionsverweis {fungsi yang terdiferensialkan secara kontinu}{}{}
dengan
\definitionsverweis {tumpuan kompak}{}{.}
Tunjukkan bahwa
\definitionsverweis {turunan}{}{}
$f'$ juga memiliki tumpuan kompak dan bahwa

\mavergleichskette
{\vergleichskette
{ \int_{\R_{\geq 0} } f' d \lambda^1
}
{ = }{ -f(0)
}
{  }{
}
{    }{
}
{   }{
}
}
{}{}{}
berlaku. Tunjukkan bahwa pernyataan ini tidak harus berlaku jika $f$ tidak memiliki tumpuan kompak.

}
{} {}

\inputaufgabe
{}
{

Misalkan $X$ suatu
\definitionsverweis {manifold diferensiabel}{}{}
dengan suatu
\definitionsverweis {basis terhitung bagi topologinya}{}{.}
Tunjukkan bahwa terdapat suatu
\definitionsverweis {struktur Riemann}{}{}
pada $X$.

}
{} {}

\inputaufgabe
{}
{

Misalkan $X$ suatu
\definitionsverweis {manifold diferensiabel}{}{}
dengan suatu
\definitionsverweis {basis terhitung bagi topologinya}{}{}
dan misalkan
\maabb {p} { E } { X
} {}
suatu
\definitionsverweis {bundel vektor diferensiabel}{}{}
di atas $X$. Tunjukkan bahwa $E$ dapat dilengkapi dengan struktur
\definitionsverweis {bundel vektor Riemann}{}{.}

}
{} {}

  \zwischenueberschrift{Soal untuk dikumpulkan}

\inputaufgabe
{3}
{

Misalkan

\mathbed {A_n} {}
 {n \in \N} {}
 {} {} {} {,}
suatu
\definitionsverweis {penghabisan kompak}{}{} dari
\definitionsverweis {ruang topologis}{}{} $X$. Dengan konvensi
\mathl{A_{-1}= \emptyset}{,}
tunjukkan bahwa relasi

\mathdisp {A_{n+1} \setminus A_n ^{o}  \subseteq A_{n+2}^{o} \setminus A_{n-1}} {  }
 berlaku.

}
{} {}

\inputaufgabe
{4}
{

Untuk
\definitionsverweis {tutupan terbuka}{}{}

\mavergleichskettedisp
{\vergleichskette
{ \R
}
{ =} { \bigcup_{n \in \Z} ]n,n+3[
}
{ } {
}
{ } {
}
{ } {
}
}
{}{}{}
berikan suatu
\definitionsverweis {partisi satuan}{}{}
kontinu yang subordinat terhadap tutupan tersebut.

}
{} {}

\inputaufgabe
{6}
{

Kita tinjau
\definitionsverweis {penghabisan kompak}{}{}

\mathbeddisp {A_n = B  \left(  0 ,n \right)} {}
 {n \in \N} {}
 {} {} {} {,}
dari $\R^2$ beserta
\definitionsverweis {tutupan terbuka}{}{}

\mathbeddisp {W_n= A_{n+1}^{o} \setminus A_{n-1}} {}
 {n \in \N} {}
 {} {} {} {,}
\zusatzklammer {dengan
\mathl{A_{-1}= \emptyset}{}} {} {.}
Carilah suatu tutupan $\R^2$ dengan
\definitionsverweis {cakram terbuka}{}{,}
yang memenuhi sifat-sifat
dari Lemma 22.9
\zusatzklammer {beserta pembuktiannya} {} {}.

}
{} {}

\inputaufgabe
{3}
{

Berikan suatu contoh
\definitionsverweis {ruang topologis}{}{,}
yang tidak memiliki
\definitionsverweis {penghabisan kompak}{}{.}

}
{} {}
